%%% data_mc_comparison_report.tex
%%% Data-MC Shower Shape Comparison with Cell-Energy Reweighting
%%% Date: 10 March 2026
%%%
%%% Compile: pdflatex data_mc_comparison_report.tex
%%%

\documentclass[a4paper,11pt]{article}

\usepackage[utf8]{inputenc}
\usepackage[T1]{fontenc}
\usepackage{graphicx}
\usepackage{amsmath}
\usepackage{amssymb}
\usepackage{booktabs}
\usepackage{hyperref}
\usepackage[margin=2.5cm]{geometry}
\usepackage{xcolor}
\usepackage{caption}
\usepackage{subcaption}
\usepackage{float}

\title{Data-MC Shower Shape Comparison\\
       with Cell-Energy Reweighting\\[1em]
       \large DAOD\_EGAM3 Ntuples}
\author{Marta Fernandes da Silva}
\date{10 March 2026}

\begin{document}

\maketitle

\begin{abstract}
This note documents the first data-vs-MC comparison of photon shower shape
variables ($R_\eta$, $R_\phi$, $w_{\eta\,2}$) using DAOD\_EGAM3 ntuples
processed through the NTupleMaker package (\texttt{dev\_Luca\_EGAM3or4} branch).
Cell-energy reweighting corrections are derived and applied following the
method of ATL-COM-PHYS-2021-640, Section~5. Two correction methods are
compared: flat shift (Method~1) and energy-dependent TProfile (Method~2).
\end{abstract}

\tableofcontents
\newpage

% =====================================================================
\section{Introduction}
% =====================================================================

Photon shower shape variables---$R_\eta$, $R_\phi$, and $w_{\eta\,2}$---are
key discriminants for photon identification in ATLAS. Mismodelling of these
variables in Monte Carlo simulation degrades the photon ID efficiency
measurement. The cell-energy reweighting method corrects individual cell
energy fractions in MC to match data, improving shower shape agreement.

This study uses DAOD\_EGAM3 (radiative $Z$-boson decay) ntuples, which
provide Layer~2 calorimeter cell energies in a $7\times11$
($\eta\times\phi$) window around the photon cluster centre.

% =====================================================================
\section{Datasets}
% =====================================================================

\begin{table}[h]
\centering
\caption{Datasets used in this study.}
\label{tab:datasets}
\begin{tabular}{lll}
\toprule
\textbf{Sample} & \textbf{Dataset} & \textbf{Grid Task ID} \\
\midrule
Data (2024) & \texttt{data24\_13p6TeV.00473235.physics\_Main} & 49030538 \\
            & \texttt{.deriv.DAOD\_EGAM3.f1444\_m2248\_p6143} & \\
\midrule
MC23e (Zee$\gamma$) & \texttt{mc23\_13p6TeV.700770.Sh\_2214\_eegamma} & 49032024 \\
                     & \texttt{.deriv.DAOD\_EGAM3.e8514\_s4369\_r16083\_p6675} & \\
\bottomrule
\end{tabular}
\end{table}

Both samples were processed with the NTupleMaker package on the
\texttt{dev\_Luca\_EGAM3or4} branch (Athena 25.0.40) using the
\texttt{-d~0} flag for DAOD input. Output ntuples contain:
\begin{itemize}
  \item \texttt{photon.cells\_7x11\_L2}: \texttt{vector<float>}, 77 cell energies
  \item \texttt{photon.ncells\_7x11\_L2}: cluster size (integer, expected = 77)
  \item \texttt{photon.reta}, \texttt{photon.unfudged\_reta}: stored (fudged/unfudged) $R_\eta$
  \item Analogous branches for $R_\phi$ and $w_{\eta\,2}$
\end{itemize}

% =====================================================================
\section{Method}
% =====================================================================

\subsection{Shower Shape Computation from Cells}

Shower shapes are computed directly from the $7\times11$ cell energy array:
\begin{align}
  R_\eta &= \frac{E(3\times7)}{E(7\times7)} \\
  R_\phi &= \frac{E(3\times3)}{E(3\times7)} \\
  w_{\eta\,2} &= \sqrt{\frac{\sum_i E_i \,\eta_i^2}{\sum_i E_i}
                 - \left(\frac{\sum_i E_i \,\eta_i}{\sum_i E_i}\right)^2}
\end{align}
where $E(m\times n)$ is the energy sum in an $m\,(\eta) \times n\,(\phi)$
sub-window centred on the cluster, and the $w_{\eta\,2}$ sum runs over a
$3\times5$ window using relative cell $\eta$ positions
$(-1, 0, +1) \times 0.025$.

Cell-computed values are validated against the stored \texttt{unfudged\_*}
branches.

\subsection{Cell-Energy Reweighting}

Following ATL-COM-PHYS-2021-640:

\paragraph{Method~1 (Flat Shift, \S5.1):}
For each cell~$k$ in each $|\eta|$ bin:
\begin{equation}
  \text{shift}_k = \langle f_{\text{data}} \rangle_k - \langle f_{\text{MC}} \rangle_k
\end{equation}
where $f_k = e_k / E_{\text{tot}}$ is the cell energy fraction. Applied as:
$f_k^{\text{corr}} = f_k^{\text{MC}} + \text{shift}_k$.

\paragraph{Method~2 (TProfile, \S5.2):}
For each cell~$k$, a TProfile captures the energy-dependent correction:
\begin{equation}
  \alpha_k(f) = \langle f_{\text{data}} - f_{\text{MC}} \rangle_{f_{\text{MC}} = f}
\end{equation}
Applied as: $f_k^{\text{corr}} = f_k^{\text{MC}} + \alpha_k(f_k^{\text{MC}})$.

In both cases, cluster energy is conserved by rescaling:
$e_k^{\text{corr}} = f_k^{\text{corr}} / \sum_j f_j^{\text{corr}} \times E_{\text{tot}}$.

% =====================================================================
\section{Initial Comparison (Before Correction)}
% =====================================================================

The initial data-vs-MC comparison shows the level of mismodelling before
any correction is applied. Distributions are normalised to unit area for
shape comparison.

Plots are generated per $|\eta|$ bin (14 bins from 0 to 2.4, excluding
the crack region $1.37 < |\eta| < 1.52$).

\textit{[Plots to be included from} \texttt{output/data\_mc\_comparison/initial\_*.pdf}\textit{]}

% Uncomment and adjust paths once plots are generated:
% \begin{figure}[H]
%   \centering
%   \includegraphics[width=0.95\textwidth]{../output/data_mc_comparison/initial_reta.pdf}
%   \caption{$R_\eta$ data-vs-MC before correction (per $|\eta|$ bin).}
% \end{figure}

% =====================================================================
\section{After Correction}
% =====================================================================

After applying the cell-energy reweighting corrections, the corrected MC
distributions are compared to data alongside the original MC.

\textit{[Plots to be included from} \texttt{output/data\_mc\_comparison/reta.pdf}\textit{,}
\texttt{rphi.pdf}\textit{,} \texttt{weta2.pdf}\textit{]}

% \begin{figure}[H]
%   \centering
%   \includegraphics[width=0.95\textwidth]{../output/data_mc_comparison/reta.pdf}
%   \caption{$R_\eta$: Data, MC, and corrected MC (Methods 1 and 2).}
% \end{figure}

% =====================================================================
\section{$\chi^2$ Summary}
% =====================================================================

The $\chi^2$ between normalised MC and data distributions quantifies the
level of agreement. Both uncorrected MC and the two correction methods
are compared.

\textit{[Table to be filled from} \texttt{validate\_data\_mc.C} \textit{output]}

\begin{table}[H]
\centering
\caption{$\chi^2$/ndf between MC and data, per $|\eta|$ bin (example format).}
\label{tab:chi2}
\begin{tabular}{c|ccc|ccc|ccc}
\toprule
$|\eta|$ bin & \multicolumn{3}{c|}{$R_\eta$} & \multicolumn{3}{c|}{$R_\phi$} & \multicolumn{3}{c}{$w_{\eta\,2}$} \\
             & MC & M1 & M2 & MC & M1 & M2 & MC & M1 & M2 \\
\midrule
0 & -- & -- & -- & -- & -- & -- & -- & -- & -- \\
\dots & \multicolumn{9}{c}{\textit{(fill after running pipeline)}} \\
\bottomrule
\end{tabular}
\end{table}

% =====================================================================
\section{Conclusions}
% =====================================================================

\textit{[To be written after reviewing the comparison plots and $\chi^2$ values.]}

\end{document}
